\chapter{Methods}\label{sec:methods}

\section{System Model}\label{sec:system-model-1}

The canonical benchmark considered in this thesis is a half-duplex
two-hop relay link comprising one source, one relay, and one destination,
with no direct source--destination path. Both hops are SISO
complex-baseband channels with independent Rayleigh fast fading. The
source uses Gray-coded QPSK, transmission is uncoded, and performance is
measured by bit-error rate (BER). The relay-processing function is the
principal experimental variable.

The canonical model is intentionally narrow. It assumes a memoryless
channel, spatially and temporally white noise, perfect receiver-side
channel-state information (CSI), and uncoded transmission. Channel
memory, intersymbol interference (ISI), imperfect channel estimation,
blind equalization, and coded block processing are introduced only in
the experiments in which those effects are explicitly studied
(Chapter~\ref{sec:unknown-channels} and
Section~\ref{sec:coded-block-df}).

\paragraph{Memorylessness and receiver information.}
In the canonical model, source symbols, fading coefficients, and noise
samples are independent across symbol indices. Consequently,
observations at indices $j\neq i$ contain no additional
population-level information about $x[i]$. With receiver-side CSI, the
pair $(y_R[i],h_1[i])$ is sufficient for symbol detection. The learned
relays considered in the canonical experiments operate on the
equalized observations defined below; unless explicitly stated
otherwise, the fading magnitude is not supplied as a separate network
input. Their learned mappings should therefore be interpreted as
regressions under the training distribution rather than as exact
CSI-conditioned posterior estimators.

\begin{figure}[htbp]
\centering
\resizebox{\textwidth}{!}{%
\begin{tikzpicture}[
  node distance=0.25cm and 0.2cm,
  box/.style={rectangle, draw, rounded corners=3pt, minimum width=1.7cm,
              minimum height=0.85cm, align=center, font=\small},
  arr/.style={->, thick, >=stealth}
]
\node[box] (src)   {Source\\Tx bits};
\node[box, right=1.2cm of src]  (mod)   {QPSK\\Mod.};
\node[box, right=1.2cm of mod]  (ch1)   {Hop 1\\Channel};
\node[box, right=1.2cm of ch1]  (relay) {Relay\\$f_\theta$ / AF / DF};
\node[box, right=1.2cm of relay](ch2)   {Hop 2\\Channel};
\node[box, right=1.2cm of ch2](demod)   {QPSK\\Demod.};
\node[box, right=1.2cm of demod]  (dst)   {Destination\\Rx bits};

\draw[arr] (src) -- (mod);
\draw[arr] (mod) -- (ch1);
\draw[arr] (ch1) -- node[above,font=\tiny]{Hop 1} (relay);
\draw[arr] (relay) -- (ch2);
\draw[arr] (ch2) -- node[above,font=\tiny]{Hop 2} (demod);
\draw[arr] (demod) -- (dst);
\end{tikzpicture}%
}
\caption[Two-hop relay system model]
{Canonical two-hop relay system. The source transmits QPSK symbols
through Hop~1, the relay applies a classical or learned processing
strategy, and the resulting symbols are transmitted through Hop~2.
Receiver-side CSI is used for one-tap zero-forcing equalization in the
canonical Rayleigh experiments.}
\label{fig:relay-system-simple}
\end{figure}

\begin{equation}\protect\phantomsection\label{eq:two-hop-system}
\text{Source}
\xrightarrow{\text{Hop 1}}
\text{Relay}
\xrightarrow{\text{Hop 2}}
\text{Destination}.
\end{equation}
\addcontentsline{equ}{listequations}
{\protect\numberline{\theequation}\quad two-hop-system}


\subsection{Modulation}\label{sec:modulation}

The canonical constellation is Gray-coded QPSK. Let
$b_0[i],b_1[i]\in\{0,1\}$ and define the corresponding bipolar axis
variables
\[
u_I[i]=1-2b_0[i],
\qquad
u_Q[i]=1-2b_1[i],
\]
where $u_I[i],u_Q[i]\in\{-1,+1\}$. The transmitted QPSK symbol is

\begin{equation}\protect\phantomsection\label{eq:qpsk-mapping}
x[i]
=
\frac{u_I[i]+j u_Q[i]}{\sqrt{2}}
=
\frac{(1-2b_0[i])+j(1-2b_1[i])}{\sqrt{2}}.
\end{equation}
\addcontentsline{equ}{listequations}
{\protect\numberline{\theequation}\quad qpsk-mapping}

The resulting constellation is
\[
\mathcal{X}_{\mathrm{QPSK}}
=
\left\{
\frac{\pm1\pm j}{\sqrt{2}}
\right\},
\]
with unit average symbol energy,
\[
E_s=\mathbb{E}[|x|^2]=1.
\]

QPSK carries $k=2$ information bits per symbol, and therefore
\[
E_b=\frac{E_s}{2}.
\]

BPSK is used primarily in the unknown-channel experiments of
Chapter~\ref{sec:unknown-channels}, where the one-dimensional alphabet
keeps exact sequence-detection benchmarks tractable. Its mapping is

\begin{equation}\protect\phantomsection\label{eq:bpsk-mapping}
x=1-2b,
\qquad
b\in\{0,1\},
\qquad
x\in\{-1,+1\}.
\end{equation}
\addcontentsline{equ}{listequations}
{\protect\numberline{\theequation}\quad bpsk-mapping}

The 16-QAM constellation appears only as one modulation-and-coding rung
in the coded link-adaptation study of
Section~\ref{sec:coded-block-df}. Higher-order constellations are not
part of the canonical relay-architecture comparison.


\subsection{Hop Model}\label{sec:hop-model}
\label{sec:two-hop-relay-model}

Communication proceeds in two orthogonal time slots because the relay
is half duplex.

\paragraph{Slot 1: source to relay.}
The relay receives

\begin{equation}\protect\phantomsection\label{eq:rayleigh-hop1}
y_R[i]
=
h_1[i]x[i]+n_1[i].
\end{equation}
\addcontentsline{equ}{listequations}
{\protect\numberline{\theequation}\quad rayleigh-hop1}

\paragraph{Slot 2: relay to destination.}
The destination receives

\begin{equation}\protect\phantomsection\label{eq:rayleigh-hop2}
y_D[i]
=
h_2[i]x_R[i]+n_2[i].
\end{equation}
\addcontentsline{equ}{listequations}
{\protect\numberline{\theequation}\quad rayleigh-hop2}

For $k\in\{1,2\}$,

\[
h_k[i]\sim\mathcal{CN}(0,1),
\qquad
n_k[i]\sim\mathcal{CN}(0,N_0),
\]

with fading and noise independent across hops and symbol indices.

Throughout this chapter, the convention

\[
\mathbb{E}[|n_k[i]|^2]=N_0
\]

is used. Consequently,

\[
\operatorname{Var}\{\Re n_k[i]\}
=
\operatorname{Var}\{\Im n_k[i]\}
=
\frac{N_0}{2}.
\]

This distinction between total complex-noise power and per-real-axis
variance is important for the QPSK and posterior-mean expressions used
below.

The nominal per-hop symbol-SNR is

\begin{equation}\protect\phantomsection\label{eq:awgn-sigma2}
\gamma_s
=
\frac{E_s}{N_0}.
\end{equation}
\addcontentsline{equ}{listequations}
{\protect\numberline{\theequation}\quad awgn-sigma2}

Since $E_s=1$ in all normalized constellations considered here,
$N_0=1/\gamma_s$. Reported SNR values therefore denote $E_s/N_0$.
For QPSK,

\[
\frac{E_b}{N_0}
=
\frac{1}{2}\frac{E_s}{N_0},
\]

so $E_b/N_0$ is $3.01$~dB below the corresponding reported
$E_s/N_0$ value.


\subsection{Receiver-Side Equalization}

The canonical Rayleigh implementation uses one-tap zero-forcing
equalization. The equalized observations are

\begin{equation}\protect\phantomsection\label{eq:awgn-hop1}
\tilde y_R[i]
=
\frac{y_R[i]}{h_1[i]}
=
x[i]+\frac{n_1[i]}{h_1[i]},
\qquad
\tilde y_D[i]
=
\frac{y_D[i]}{h_2[i]}
=
x_R[i]+\frac{n_2[i]}{h_2[i]}.
\end{equation}
\addcontentsline{equ}{listequations}
{\protect\numberline{\theequation}\quad awgn-hop1}

Conditioned on $h_k[i]$, the zero-forced noise has total complex
variance

\[
\mathbb{E}
\left[
\left|
\frac{n_k[i]}{h_k[i]}
\right|^2
\middle|
h_k[i]
\right]
=
\frac{N_0}{|h_k[i]|^2},
\]

and variance per real component

\[
\frac{N_0}{2|h_k[i]|^2}.
\]

The instantaneous post-equalization SNR is therefore

\[
\gamma_k[i]
=
|h_k[i]|^2\gamma_s.
\]

Zero forcing should not be confused with phase compensation. Phase
compensation forms

\[
y_k[i]\frac{h_k^*[i]}{|h_k[i]|}
=
|h_k[i]|x[i]
+
n_k[i]e^{-j\arg h_k[i]},
\]

which removes the channel phase but retains its magnitude. The two
statistics satisfy

\[
y_k[i]\frac{h_k^*[i]}{|h_k[i]|}
=
|h_k[i]|
\frac{y_k[i]}{h_k[i]},
\]

for $|h_k[i]|>0$. They therefore induce the same ideal hard-decision
regions, but they are not equivalent soft observations. In particular,
they differ in amplitude, conditional noise variance, relay power
normalization, and posterior metrics. The canonical experiments use
the zero-forced statistic defined in
Equation~\eqref{eq:awgn-hop1}.

The AWGN special case is obtained by setting $h_k[i]=1$, giving

\[
y_R[i]=x[i]+n_1[i],
\qquad
y_D[i]=x_R[i]+n_2[i].
\]

This special case is used where analytical AWGN expressions are
required; no separate AWGN-only canonical relay study is implied.


\subsection{Evaluated Configurations}
\label{sec:evaluated-configurations}

Table~\ref{tbl:configurations} summarizes the principal
constellation/channel combinations.

\begin{longtable}[]{@{}lllp{0.42\linewidth}@{}}
\caption[The evaluated configurations]
{Primary constellation/channel configurations used in the thesis.}
\label{tbl:configurations}\tabularnewline
\toprule
Constellation & Channel & Role & Purpose \\
\midrule
\endfirsthead

\toprule
Constellation & Channel & Role & Purpose \\
\midrule
\endhead

\bottomrule
\endlastfoot

QPSK (complex)
&
Rayleigh (complex)
&
Canonical
&
Primary relay-architecture comparison. QPSK exercises both dimensions
of the complex channel while retaining a binary decision on each axis.
\\

BPSK (real)
&
Unknown three-tap ISI
&
Extension
&
Used in Chapter~\ref{sec:unknown-channels} to study channel memory while
keeping exact sequence-detection benchmarks computationally compact.
\\

16-QAM (complex)
&
Rayleigh (complex)
&
Coded MCS rung
&
Used only in Section~\ref{sec:coded-block-df} as part of the
modulation-and-coding grid.
\\

\end{longtable}

Coherent BPSK over complex Rayleigh fading is theoretically valid and
is not excluded for physical reasons. QPSK is used canonically as an
experimental design choice because it exercises a genuinely
two-dimensional complex constellation while retaining independent
binary decisions on the in-phase and quadrature axes.


\subsection{Relay Processing}\label{sec:relay-processing-methods}

The relay-processing function is written generally as

\[
x_R[i]
=
f_\theta(\mathcal{Y}_{R,i}),
\]

where $\mathcal{Y}_{R,i}$ denotes the observation supplied to the relay.
For AF and symbol-wise DF,

\[
\mathcal{Y}_{R,i}=\tilde y_R[i].
\]

For learned relays, a finite observation window may be used.

Because the canonical QPSK symbols have axis amplitudes
$\pm1/\sqrt{2}$, it is convenient to define normalized real-valued
axis observations

\[
r_I[i]
=
\sqrt{2}\Re\{\tilde y_R[i]\},
\qquad
r_Q[i]
=
\sqrt{2}\Im\{\tilde y_R[i]\}.
\]

The corresponding clean targets are

\[
u_I[i],u_Q[i]\in\{-1,+1\}.
\]

A shared real-valued relay function can then be applied independently
to the two axes:

\begin{equation}\protect\phantomsection\label{eq:iq-splitting}
\hat x_R[i]
=
\frac{
f_\theta(\mathbf r_{I,i})
+
j f_\theta(\mathbf r_{Q,i})
}{\sqrt{2}}.
\end{equation}
\addcontentsline{equ}{listequations}
{\protect\numberline{\theequation}\quad iq-splitting}

\textbf{Implementation verification note.}
The final archived release should be checked to confirm that the
implementation explicitly applies the $\sqrt{2}$ axis normalization
shown above and that the same parameter object is shared between the
I and Q evaluations. If the implementation instead trains directly on
targets $\pm1/\sqrt{2}$, Equation~\eqref{eq:iq-splitting} should be
modified accordingly rather than relying on subsequent power
normalization to absorb the scale difference.


\paragraph{Posterior-mean interpretation.}

For a normalized binary axis symbol $u\in\{-1,+1\}$ observed through

\[
r=u+v,
\qquad
v\sim\mathcal{N}(0,\sigma_r^2),
\]

with equal symbol priors, the MSE-optimal estimator is the posterior
mean

\[
\mathbb{E}[u\mid r]
=
\tanh\left(\frac{r}{\sigma_r^2}\right).
\]

Under the canonical zero-forced Rayleigh model,

\[
\sigma_r^2
=
\frac{N_0}{|h|^2},
\]

after the $\sqrt{2}$ axis normalization above. Hence, when $h$ is
known,

\[
\mathbb{E}[u\mid r,h]
=
\tanh\left(
\frac{|h|^2r}{N_0}
\right).
\]

This expression motivates the use of a bounded nonlinear regression
function, but it should not be interpreted as the exact function
implemented by a network that does not receive $|h|$ explicitly. Such
a network instead approximates the marginal regression

\[
\mathbb{E}[u\mid r]
\]

under the fading and SNR distribution represented in its training
data.


\paragraph{Relay power normalization.}

Relay outputs are normalized to a common target average power before
Hop~2. For a block of $B$ relay outputs, define

\[
\widehat P_R
=
\frac{1}{B}
\sum_{i=1}^{B}
|x_R[i]|^2.
\]

The normalized relay output is

\begin{equation}\protect\phantomsection\label{eq:power-normalization}
x_R[i]
\leftarrow
x_R[i]
\sqrt{
\frac{P_{\mathrm{target}}}
{\widehat P_R}
}.
\end{equation}
\addcontentsline{equ}{listequations}
{\protect\numberline{\theequation}\quad power-normalization}

\textbf{Implementation verification note.}
The final release should be checked to confirm whether
$\widehat P_R$ is computed per transmitted block, minibatch, trial, or
another scope. The text above should be adjusted to match the exact
implementation.


\begin{remark}[Realizability of the symmetric window]
\label{rem:window-causality}

The canonical channel is memoryless. Therefore, under the assumed
independence of source symbols, fading, and noise,

\[
p(x[i]\mid y_R[i-w:i+w],h_1[i-w:i+w])
=
p(x[i]\mid y_R[i],h_1[i]).
\]

Neighboring observations consequently provide no additional
population-level information about $x[i]$.

Some learned relays nevertheless receive a symmetric window

\[
\mathbf r_i
=
\bigl(
r[i-w],\ldots,r[i],\ldots,r[i+w]
\bigr).
\]

This is an architectural convention rather than a model of channel
memory. It also permits related windowed architectures to be reused in
the ISI experiments of Chapter~\ref{sec:unknown-channels}, where
neighboring observations do contain useful information.

A symmetric window is not streaming-causal because it requires $w$
future observations. It is realizable in the half-duplex block
protocol if the relay buffers the listening-phase block before
transmission, at the cost of a fixed look-ahead latency.

The theoretical redundancy of neighboring samples should not be
confused with an empirical invariance claim. Finite trained networks
may react to redundant coordinates because of finite-sample effects,
optimization, padding, or normalization. No causal-window ablation is
reported in the canonical study; therefore, the results do not
establish that replacing the symmetric window by a causal or
single-sample input would leave the measured BER unchanged.

\end{remark}


\section{Relay Strategies}\label{sec:relay-strategies}

Nine relay strategies were implemented, comprising two classical
baselines and seven learned strategies. They span feedforward,
hybrid, variational, adversarial, attention-based, and state-space
architectures.

The purpose of this set is not to establish an exhaustive taxonomy of
relay processing. Rather, it provides representative architectures
with substantially different parameter counts and inductive
structures under a common end-to-end relay experiment.

\subsection{Selection and Classification of Relay Strategies}

The strategies are grouped as follows:

\begin{enumerate}
    \item classical baselines: AF and symbol-wise DF;
    \item supervised feedforward relays: MLP and Hybrid;
    \item generative-objective relays: VAE and cGAN;
    \item sequence architectures: Transformer, Mamba-S6, and
          Mamba-2 SSD.
\end{enumerate}


\subsection{Amplify-and-Forward}

Amplify-and-forward (AF) applies no learned denoising or hard
regeneration. The relay forwards a scaled version of its received
observation,

\[
x_R[i]
=
G\tilde y_R[i],
\]

where $G$ is selected according to the relay power-normalization
convention.

AF therefore preserves both the desired signal and the first-hop
noise. It serves as the minimal-processing classical baseline.


\subsection{Symbol-Wise Decode-and-Forward}

The canonical DF relay performs hard QPSK demodulation on the
equalized first-hop observation and immediately remodulates the
detected bits. Processing is performed independently for each symbol,
with no outer channel code.

For QPSK,

\[
\hat b_0[i]
=
\mathbb{1}
\left[
\Re\{\tilde y_R[i]\}<0
\right],
\]

\[
\hat b_1[i]
=
\mathbb{1}
\left[
\Im\{\tilde y_R[i]\}<0
\right],
\]

followed by QPSK remodulation.

DF is therefore the hard-regeneration baseline evaluated in the
canonical uncoded comparison.


\begin{remark}[Terminology: symbol-wise DF versus information-theoretic DF]
\label{rem:df-terminology}

The term ``decode-and-forward'' is used in two related but distinct
senses.

In information-theoretic relay-channel models
\cite{CoverGamal1979RelayChannel,ElGamalKim2011NetworkIT}, DF normally
denotes a block strategy in which the relay decodes a coded message
block and subsequently re-encodes and forwards the recovered message.

The canonical experiments in this thesis are uncoded. Consequently,
``DF'' in the canonical comparison denotes symbol-wise hard slicing
and remodulation rather than coded block decoding. It is the
appropriate hard-regeneration baseline for the uncoded BER experiment,
but its performance should not be interpreted as a bound on coded
block-DF.

A practical finite-length coded block-DF baseline is evaluated
separately in Section~\ref{sec:coded-block-df}.

\end{remark}


\subsection{Minimal MLP Relay}

The minimal learned relay is a two-layer feedforward multilayer
perceptron. For one normalized QPSK axis, the input is a five-sample
window,

\[
\mathbf r_i
=
[r[i-2],r[i-1],r[i],r[i+1],r[i+2]]^\mathsf{T}.
\]

The network is

\begin{equation}\label{eq:mlp-relay}
f_\theta(\mathbf r)
=
\tanh
\left(
\mathbf W_2
\operatorname{ReLU}
(\mathbf W_1\mathbf r+\mathbf b_1)
+b_2
\right).
\end{equation}

With five inputs and 24 hidden units, the parameter count is

\[
(5\times24+24)+(24\times1+1)=169.
\]

The same network is applied to the normalized I and Q streams.

The bounded tanh output is compatible with the normalized binary
targets $\{-1,+1\}$. The architecture can be interpreted as a
low-capacity nonlinear regression model for estimating a clean binary
axis symbol from a noisy local observation.

Training uses MSE loss and multi-SNR synthetic training data generated
at 5, 10, and 15~dB.

The reported implementation uses 25,000 supervised training samples,
100 epochs, and learning rate $\eta=0.01$.

\textbf{Implementation verification note.}
The final archived release should be used to verify the exact optimizer,
batching convention, initialization, and whether the 25,000-sample
count denotes total examples or examples per training SNR.


\subsection{Hybrid MLP--DF Relay}

The Hybrid relay combines the trained MLP with symbol-wise DF. A
switching rule selects between the two relay functions according to
the operating SNR.

The switching threshold is determined empirically from the training or
validation comparison between the MLP and DF rather than introduced as
an additional learned neural parameter. The Hybrid therefore has the
same neural parameter count as the MLP.

This strategy tests whether a simple operating-condition-dependent
selection rule can exploit complementary behavior of soft learned
regression and hard symbol regeneration. Claims about which method is
superior in a particular SNR regime are treated as experimental
results rather than assumptions of the method.


\subsection{Variational Autoencoder Relay}

The VAE relay uses a conditional reconstruction architecture with an
encoder that maps the noisy observation window to a latent
distribution and a decoder that reconstructs the clean target.

The reported architecture is

\[
7
\rightarrow
32
\rightarrow
16
\rightarrow
(\boldsymbol\mu,\boldsymbol\sigma^2),
\]

with latent dimension 8, followed by

\[
8
\rightarrow
16
\rightarrow
32
\rightarrow
1.
\]

The reported parameter count is 1,777.

Training uses a $\beta$-VAE objective,

\[
\mathcal L_{\mathrm{VAE}}
=
\mathcal L_{\mathrm{rec}}
+
\beta
D_{\mathrm{KL}}
\left(
q_\phi(\mathbf z\mid\mathbf r)
\Vert
p(\mathbf z)
\right),
\]

with $\beta=0.1$.

The relatively small KL weight places greater emphasis on
reconstruction than on latent regularization. This is a design choice;
the thesis does not claim that $\beta=0.1$ is optimal.


\subsection{Conditional GAN Relay}

The cGAN relay uses a generator conditioned on the noisy observation
and a critic trained with a Wasserstein objective and gradient penalty.

The reported generator architecture is

\[
(7+8)
\rightarrow
32
\rightarrow
32
\rightarrow
16
\rightarrow
1,
\]

and the critic is

\[
(1+7)
\rightarrow
32
\rightarrow
16
\rightarrow
1.
\]

The reported total parameter count is 2,946.

The training objective combines the adversarial term with an L1
reconstruction term,

\[
\mathcal L_G
=
\mathcal L_{\mathrm{adv}}
+
\lambda_{\mathrm{L1}}
\mathcal L_{\mathrm{L1}},
\]

with

\[
\lambda_{\mathrm{L1}}=100,
\qquad
\lambda_{\mathrm{GP}}=10.
\]

The adversarial objective therefore supplements a strongly weighted
conditional reconstruction objective. The model is described as
generative because of its training formulation; its role in the relay
experiment is conditional signal reconstruction.


\subsection{Transformer Relay}

The Transformer provides an attention-based sequence-model baseline.
It operates on an 11-sample observation window.

The reported configuration is

\[
d_{\mathrm{model}}=32,
\qquad
n_{\mathrm{heads}}=4,
\qquad
n_{\mathrm{layers}}=2,
\qquad
d_{\mathrm{ff}}=128.
\]

The reported parameter count is 17,697.

The model is trained for 100 epochs using Adam with learning rate
$10^{-3}$.

The Transformer is included as a widely used attention-based sequence
architecture. Its inclusion does not imply that attention is expected
a priori to be advantageous on the memoryless canonical channel.


\subsection{Mamba-S6 Relay}

The Mamba-S6 relay is a selective state-space sequence model with
input-dependent state-space parameters. The reported configuration is

\[
d_{\mathrm{model}}=32,
\qquad
d_{\mathrm{state}}=16,
\]

with two Mamba blocks and residual connections.

The reported parameter count is 24,001.

Each block contains normalization, input projection, local convolution
and nonlinear processing, a selective state-space operation, gating,
output projection, and a residual connection.

The state dimension $d_{\mathrm{state}}=16$ specifies the internal
state size. It should not be interpreted directly as the order of a
classical linear adaptive filter because the Mamba transition and
selection mechanisms are input dependent and nonlinear at the block
level.


\subsection{Mamba-2 SSD Relay}

The Mamba-2 relay uses the Structured State Space Duality (SSD)
formulation. The reported configuration is

\[
d_{\mathrm{model}}=32,
\qquad
d_{\mathrm{state}}=16,
\qquad
\text{chunk size}=8,
\]

with two Mamba-2 blocks.

The reported parameter count is 26,179.

Training uses Adam with learning rate $10^{-3}$ and gradient clipping

\[
\|\nabla\|_2\leq1.
\]

The SSD formulation provides a structured implementation of the
state-space computation. The chunk size is an implementation
hyperparameter controlling the organization of the computation rather
than a property of the physical relay channel.


\section{Simulation Framework}\label{sec:simulation-framework}

\subsection{Monte Carlo BER Estimation}
\label{sec:monte-carlo-ber-estimation}

For the canonical relay comparison, BER is estimated using independent
Monte Carlo trials at each SNR point.

The nominal canonical simulation budget is

\begin{itemize}
    \item bits per trial: $N=10{,}000$;
    \item trials per SNR: $M=10$;
    \item total transmitted bits per SNR point: $MN=100{,}000$;
    \item SNR range: 0--20~dB in 2~dB increments.
\end{itemize}

For trial $m$, the BER estimate is

\begin{equation}\protect\phantomsection\label{eq:ber-estimate}
\hat P_e^{(m)}
=
\frac{1}{N}
\sum_{i=1}^{N}
\mathbb{1}
\left(
b_i^{(m)}
\neq
\hat b_i^{(m)}
\right).
\end{equation}
\addcontentsline{equ}{listequations}
{\protect\numberline{\theequation}\quad ber-estimate}

The reported Monte Carlo mean is

\[
\bar P_e
=
\frac{1}{M}
\sum_{m=1}^{M}
\hat P_e^{(m)}.
\]


\subsection{Uncertainty Quantification}

When uncertainty is summarized across independent trials, the
trial-level standard error is

\[
\operatorname{SE}(\bar P_e)
=
\frac{s}{\sqrt M},
\]

where

\[
s^2
=
\frac{1}{M-1}
\sum_{m=1}^{M}
\left(
\hat P_e^{(m)}-\bar P_e
\right)^2.
\]

A descriptive 95\% Student-$t$ interval is therefore

\begin{equation}\protect\phantomsection\label{eq:confidence-interval}
\bar P_e
\pm
t_{0.975,M-1}
\frac{s}{\sqrt M}.
\end{equation}
\addcontentsline{equ}{listequations}
{\protect\numberline{\theequation}\quad confidence-interval}

For $M=10$,

\[
t_{0.975,9}\approx2.262.
\]

This interval describes variation across Monte Carlo trials. At very
low BER, however, trial estimates become discrete and may contain many
zeros, so the normal/$t$ approximation becomes less informative.

A BER value of zero should therefore be interpreted as ``no errors
observed within the finite simulation budget,'' not as evidence that
the true BER is zero. The increment associated with one observed error
in a pooled $MN$-bit experiment is

\[
\frac{1}{MN},
\]

but this is a numerical resolution, not a statistical detection
threshold. When no errors are observed, an upper confidence bound
should be reported or discussed where quantitative low-BER inference
is required.


\subsection{Per-Experiment Simulation Budgets}

The $10\times10{,}000$-bit budget applies to the canonical relay
comparison and the associated parameter-normalization study.

Experiments involving unknown or randomly redrawn channel families use
larger budgets as specified in their respective sections. In
particular, increasing the number of bits within one fixed channel
realization reduces conditional Monte Carlo noise, whereas increasing
the number of independent trials is required to average over a family
of randomly redrawn channel realizations.

For this reason, experiments in which the impairment itself is redrawn
between blocks use larger trial counts where specified in
Chapter~\ref{sec:unknown-channels}.


\subsection{Statistical Significance Testing}
\label{sec:statistical-significance-testing}

Pairwise differences between relay strategies are assessed using the
Wilcoxon signed-rank test when paired trial observations are available.

For methods $A$ and $B$, define

\[
d_m
=
\hat P_{e,A}^{(m)}
-
\hat P_{e,B}^{(m)}.
\]

The hypotheses are

\begin{equation}\protect\phantomsection\label{eq:wilcoxon-test}
H_0:
\operatorname{median}(d_m)=0
\qquad
\text{versus}
\qquad
H_1:
\operatorname{median}(d_m)\neq0.
\end{equation}
\addcontentsline{equ}{listequations}
{\protect\numberline{\theequation}\quad wilcoxon-test}

The paired interpretation requires methods being compared to use the
same underlying trial realization---including the same transmitted
bits and, where applicable, the same fading and noise realization.
This common-random-number design should be verified against the
archived experiment scripts for every experiment in which the
signed-rank test is reported.

With ten non-zero untied paired differences, the smallest attainable
exact two-sided Wilcoxon $p$-value is

\[
\frac{2}{2^{10}}
=
0.001953125.
\]

Zero differences reduce the effective sample size and therefore change
the attainable exact $p$-values.

Unless otherwise stated, $\alpha=0.05$ is interpreted per comparison.
No global family-wise or false-discovery-rate correction is implied.
Consequently, collections of many pairwise tests should be interpreted
as exploratory unless a multiple-comparison correction is explicitly
reported.


\subsection{Training Protocol}\label{sec:training-protocol}

Learned relays are trained offline before the BER evaluation sweep.
The same trained weights are then reused across the evaluated SNR
points without further parameter updates.

Training data are generated synthetically at

\[
\mathrm{SNR}
\in
\{5,10,15\}\ \mathrm{dB}.
\]

The purpose of multi-SNR training is to expose each model to several
noise levels rather than fitting a separate model at every evaluation
SNR.

The evaluation grid extends beyond the training grid. Performance at
SNR values between the training points therefore involves
interpolation with respect to the training distribution, whereas
evaluation below 5~dB or above 15~dB involves some degree of
distributional extrapolation. Whether this produces a measurable BER
penalty is an experimental question and is not assumed in the method.

AF and DF contain no trainable neural parameters.

The learned strategies use model-specific objectives:

\begin{itemize}
    \item MLP and Hybrid MLP component: supervised regression;
    \item VAE: reconstruction plus KL regularization;
    \item cGAN: adversarial plus reconstruction objective;
    \item Transformer, Mamba-S6, and Mamba-2 SSD:
          supervised sequence-model training.
\end{itemize}

The final archived implementation should be treated as authoritative
for exact sample counts, batch sizes, optimizer settings, initialization
rules, and checkpoint behavior where these differ between model
families.


\subsection{Reproducibility and Weight Management}
\label{sec:reproducibility-and-weight-management}

Experiments use controlled pseudorandom seeds for source generation,
channel realizations, and noise generation. Trained model parameters
are checkpointed so that BER evaluation can be repeated without
retraining.

For reproducibility, the archived experimental release should record

\begin{itemize}
    \item the exact source-code revision;
    \item software and dependency versions;
    \item model hyperparameters;
    \item random seeds;
    \item training and evaluation budgets;
    \item parameter counts;
    \item result-file hashes or equivalent provenance information.
\end{itemize}

The simulation framework and experiment scripts are maintained in the
\texttt{relaynet} project repository.

Rather than embedding a mutable count of passing automated tests in the
thesis prose, the archived release should be regarded as the
authoritative record of the test-suite result associated with the
reported experiments. This avoids inconsistencies when additional
tests are added after an earlier thesis draft is generated.


\section{Normalized Parameter Comparison}
\label{sec:normalized-parameter-comparison}

The canonical learned architectures differ substantially in trainable
parameter count. Consequently, a direct comparison confounds
architecture with model capacity.

To reduce this confound, a secondary experiment scales all seven
learned relay families to approximately 3,000 trainable parameters.

The purpose of this experiment is limited: it controls the scalar
number of trainable parameters to first order. It does not isolate
architectural inductive bias because reaching a common parameter
budget requires architecture-specific changes to width, depth, latent
dimension, state dimension, number of attention heads, and related
hyperparameters.

\begin{longtable}[]{@{}lll@{}}
\caption[Normalized 3K-parameter model configurations]
{Approximately 3K-parameter configurations used in the
capacity-normalized comparison.}
\label{tbl:table28}\tabularnewline

\toprule
Model & Parameters & Scaling method \\
\midrule
\endfirsthead

\toprule
Model & Parameters & Scaling method \\
\midrule
\endhead

\bottomrule
\endlastfoot

MLP-3K
&
3,004
&
Window $5\rightarrow11$; hidden width $24\rightarrow231$
\\

Hybrid-3K
&
3,004
&
Same MLP scaling plus DF switching
\\

VAE-3K
&
3,037
&
Window $7\rightarrow11$; latent/hidden dimensions adjusted
\\

cGAN-3K
&
3,004
&
Window $7\rightarrow11$; hidden dimensions adjusted
\\

Transformer-3K
&
3,007
&
$d_{\mathrm{model}}:32\rightarrow18$;
heads $4\rightarrow2$;
layers $2\rightarrow1$
\\

Mamba-S6-3K
&
3,027
&
$d_{\mathrm{model}}:32\rightarrow16$;
$d_{\mathrm{state}}:16\rightarrow6$;
layers $2\rightarrow1$
\\

Mamba-2-SSD-3K
&
3,004
&
$d_{\mathrm{model}}:32\rightarrow15$;
$d_{\mathrm{state}}:16\rightarrow6$;
layers $2\rightarrow1$
\\

\end{longtable}

All normalized configurations use an 11-sample input window. This
removes the large window-length differences present in the original
feedforward, generative, and sequence configurations within this
specific comparison.

All learnable parameters are counted, including biases, normalization
gain/bias parameters, and learned projection matrices. Fixed buffers
and non-trainable constants are excluded.

The parameter count is verified programmatically using the equivalent
of

\begin{verbatim}
sum(
    p.numel()
    for p in model.parameters()
    if p.requires_grad
)
\end{verbatim}

The normalized experiment therefore removes large differences in
trainable-parameter count and equalizes nominal input-window length.
It does \emph{not} isolate architecture. Width, depth, latent/state
dimension, attention-head count, computational cost, and optimization
behavior remain architecture dependent.

Accordingly, if performance differences remain after normalization,
the defensible conclusion is that parameter count alone does not
explain those differences. The experiment does not by itself establish
that architectural inductive bias is the unique cause.


\section{QPSK Processing}\label{sec:modulation-schemes}

The canonical QPSK mapping was defined in
Equation~\eqref{eq:qpsk-mapping}. This section records the corresponding
hard decisions and explains the axis-separable processing used by the
real-valued learned relays.

\subsection{Gray-Coded QPSK}
\label{sec:qpsk-gray-coded-quadrature-phase-shift-keying}

The four QPSK symbols are

\[
\left\{
\frac{+1+j}{\sqrt2},
\frac{+1-j}{\sqrt2},
\frac{-1-j}{\sqrt2},
\frac{-1+j}{\sqrt2}
\right\}.
\]

The corresponding mapping is shown in
Table~\ref{tbl:table29}.

\begin{longtable}[]{@{}lll@{}}
\caption[QPSK Gray-coded symbol mapping]
{Gray-coded QPSK bit-to-symbol mapping.}
\label{tbl:table29}\tabularnewline

\toprule
Bit pair $(b_0,b_1)$ & Symbol & Quadrant \\
\midrule
\endfirsthead

\toprule
Bit pair $(b_0,b_1)$ & Symbol & Quadrant \\
\midrule
\endhead

\bottomrule
\endlastfoot

00 & $(+1+j)/\sqrt2$ & I \\
01 & $(+1-j)/\sqrt2$ & IV \\
11 & $(-1-j)/\sqrt2$ & III \\
10 & $(-1+j)/\sqrt2$ & II \\

\end{longtable}

Hard demodulation is performed independently on the two axes:

\[
\hat b_0
=
\mathbb{1}
\left[
\Re\{\hat x\}<0
\right],
\qquad
\hat b_1
=
\mathbb{1}
\left[
\Im\{\hat x\}<0
\right].
\]

For uncoded Gray-coded QPSK over AWGN, the BER at a given
$E_b/N_0$ equals the coherent BPSK BER,

\begin{equation}\protect\phantomsection\label{eq:qpsk-ber}
P_b^{\mathrm{QPSK}}
=
Q
\left(
\sqrt{\frac{2E_b}{N_0}}
\right).
\end{equation}
\addcontentsline{equ}{listequations}
{\protect\numberline{\theequation}\quad qpsk-ber}

This equality is stated in terms of $E_b/N_0$. Since the simulation
SNR convention in this thesis is $E_s/N_0$, the corresponding QPSK
$E_b/N_0$ value is 3.01~dB lower.


\subsection{I/Q Splitting for Learned Relays}
\label{sec:iq-splitting-for-ai-relay-processing-of-complex-constellations}

For QPSK, the clean normalized axis variables $u_I$ and $u_Q$ are
independent binary variables. After coherent equalization, the
corresponding real noise components are also independent conditioned
on the channel coefficient.

This motivates the axis-separable relay representation

\[
\hat x_R
=
\frac{
f_\theta(\mathbf r_I)
+
j f_\theta(\mathbf r_Q)
}{\sqrt2}.
\]

For the canonical QPSK problem, this decomposition is particularly
natural because each axis contains one binary decision.

It should nevertheless be described as a restriction of the relay
hypothesis class rather than as a general equivalence between
real-valued axis processing and arbitrary complex-valued processing.
A joint complex-valued relay could, in principle, represent mappings
that couple the two axes.

Under the canonical channel assumptions and binary-per-axis QPSK
decision problem, no such coupling is required by the ideal coherent
symbol detector. This does not imply that the same decomposition is
lossless for arbitrary learned objectives, channel models, or
higher-order constellations.

For constellations such as 16-QAM, each axis contains more than two
amplitude levels. A relay with a single tanh output trained only on
binary $\{-1,+1\}$ targets is therefore not directly transferable
without changing the target representation, retraining the model, or
using a joint multidimensional output.


\subsection{Constellation Diagram Summary}
\label{sec:constellation-diagram-summary}

Figure~\ref{fig:fig8} illustrates the modulation alphabets appearing
in the thesis. QPSK is the canonical relay-comparison constellation,
BPSK is used in the principal unknown-ISI experiments, and 16-QAM
appears only in the coded modulation-and-coding study. Other
constellations shown in the figure are included for orientation rather
than as evaluated canonical relay configurations.

\begin{figure}[htbp]
\centering
\pandocbounded{
\includegraphics[
width=\linewidth,
height=0.45\textheight,
keepaspectratio
]{results/modulation/constellation_diagrams.png}
}
\caption[Constellation diagrams]
{Constellation diagrams used for orientation. BPSK contains two
real-valued symbols, QPSK contains four unit-energy Gray-coded symbols,
and 16-QAM uses a denser rectangular grid. QPSK is the canonical
constellation used in the relay-architecture comparison; BPSK and
16-QAM are used only in the extension studies identified in the text.}
\label{fig:fig8}
\end{figure}


\section{Methodological Scope and Verification Notes}

The canonical methodology is intentionally separated from the
extension studies.

The conclusions drawn from the canonical experiments therefore apply
to the following setting:

\begin{itemize}
    \item half-duplex two-hop SISO relaying;
    \item no direct source--destination link;
    \item independent complex Rayleigh fast fading on both hops;
    \item perfect receiver-side CSI;
    \item one-tap zero-forcing equalization;
    \item Gray-coded QPSK;
    \item uncoded BER;
    \item offline-trained learned relays;
    \item common relay-output power normalization.
\end{itemize}

Channel memory, unknown channels, finite-pilot estimation, blind
operation, coded block-DF, and modulation-and-coding adaptation are
deliberate departures from this model and are defined where they are
evaluated.

Before the final thesis release, the following implementation details
should be verified directly against the exact archived code revision
that generated the reported result files:

\begin{enumerate}
    \item confirm that every canonical Rayleigh result uses the
          zero-forcing code path;
    \item confirm the exact QPSK I/Q scaling used at learned-relay
          inputs and outputs;
    \item confirm whether one parameter object is shared between I and
          Q processing;
    \item confirm the scope of relay-output power normalization;
    \item confirm model-specific training-sample counts, batch sizes,
          optimizers, and initialization rules;
    \item confirm the exact cGAN inclusion/exclusion policy in each
          reported comparison;
    \item confirm that paired statistical tests use common underlying
          Monte Carlo realizations;
    \item confirm the treatment of zero observed BER values in stored
          results and plots;
    \item record the exact repository revision associated with the
          final thesis results.
\end{enumerate}

These checks concern reproducibility and implementation provenance;
they do not alter the mathematical definition of the canonical system
given above.